\runninghead{The Three Interpretations Compared}
\section{The Three Interpretations Compared}
\label{sec:comparison}

\begin{comparisontable}{}{Where the body is}{Not coming down}{Words that do not pass}
Question &
What is the end a gathering to? &
To whom are the commands of flight spoken? &
What makes a promise certain when its day is unknown? \\
Starting words &
\latin{Ubicúmque fúerit corpus, illic congregabúntur et áquilæ}, Mt 24:28;
\latin{congregábunt eléctos eius}, v.~31 &
\latin{qui in tecto, non descéndat}, Mt 24:17 &
\latin{verba autem mea non præteríbunt}, Mt 24:35 \\
Principal authors &
Jerome, Hilary, Augustine &
Hilary, Jerome &
Chrysostom, Gregory, Augustine \\
Its answer &
To Christ crucified and risen: the Passion as read, the place of the
Passion, the body in heaven &
To every hearer as well as to those in Jerusalem: who has gone up must not
go back down &
The speaker: the past fulfilled, and the future as certain with God as the
past \\
In the Introit &
The captivity brought back from all places &
A letter that warned the exiles against their prophets &
God's thoughts, known to him, and prayer that he will hear \\
In the Epistle &
Delivered from darkness and translated into the kingdom (1:13) &
Walking worthy in patience and longsuffering (1:10--11) &
Thanks for what is given, petition for what is not yet (1:9--12) \\
Weight among the senses &
Allegorical and anagogical &
Moral &
Allegorical and moral \\
\end{comparisontable}

The three interpretations read the same Gospel with largely the same Fathers.
All three take its last part of the Lord's visible coming, and the second sets
Chrysostom's literal sense of the flight first. They differ in the verse from
which each begins and in what each asks of the Gospel: where the end leads,
how to live before it, and why it can be trusted.

The same texts serve each question differently. The Introit's return from all
places is, in the first, the beginning of the gathering of the scattered. The
second draws instead on the letter the antiphon comes from, in which Jeremias
warned the exiles against their prophets and diviners, as the Gospel warns
against false christs. In the third it is a promise whose keeping God alone
foresees and for which he still bids his people pray. The Gradual's
\latin{Liberásti} is the thanks of those brought back (Rupert), the present
deliverance of those whom afflicters cannot harm in spirit (Aquinas), and the
past tense of a certain future (Augustine). From the Secret the first takes
the hearts of all turned to one God, the second the passage from earthly
desires, and the third the desires for what does not pass.

The second and third rest on Fathers who disagree about method. Hilary's step
from the letter of the housetop to the soul, which the second follows, is the
step Chrysostom does not take; yet the third interpretation is built chiefly
on Chrysostom. Neither leaves the other Father out: the second sets
Chrysostom's war first, and the third uses Hilary on the generation that does
not pass and on the words that remain. On the fig tree it does not follow him,
since Hilary's branch is Antichrist and Chrysostom's and Gregory's is a pledge
of summer.
